% TOP_PROBLEM: {"id": 3168, "title": "Hamiltonian cycles in line graphs", "category_id": 3, "category": {"id": 3, "name": "graph_theory", "display_name": "Graph Theory", "description": "Problems involving graphs, networks, and their properties.", "slug": "graph-theory", "order_index": 3, "created_at": "2026-07-31T15:26:25.671Z"}, "status": "open", "problem_number": "OPG-485", "set_id": 12}
% FIRST_POSTED: 2026-10-01
% ATTEMPT_STATUS: unresolved
% UnsolvedMath ID 3168; source catalogue code OPG-485
% !TeX program = lualatex
% GPT-6 Astra Ultra research attempt; independent partial work, 2026-10-01.
\documentclass[11pt,a4paper]{article}
\usepackage[margin=25mm]{geometry}
\usepackage{fontspec}
\setmainfont{lmroman10-regular.otf}[BoldFont=lmroman10-bold.otf,
 ItalicFont=lmroman10-italic.otf,BoldItalicFont=lmroman10-bolditalic.otf]
\usepackage{amsmath,amssymb,mathtools}
\usepackage{unicode-math}
\setmathfont{latinmodern-math.otf}
\AtBeginDocument{\let\setminus\smallsetminus}
\usepackage{xurl,hyperref}
\hypersetup{colorlinks=true,urlcolor=blue}
\setcounter{secnumdepth}{0}
\setlength{\parindent}{0pt}
\setlength{\parskip}{5pt}
\setlength{\emergencystretch}{3em}
\begin{document}
\section*{Hamiltonian cycles in line graphs}\label{3168}
{\small UnsolvedMath ID 3168 · OPG-485 · Graph Theory · OpenGarden}
\subsection{Definitions and mathematical statement}
For a finite simple undirected graph $H$, the line graph
$L(H)$ has vertex set $E(H)$. Distinct edges of $H$ become
adjacent vertices of $L(H)$ exactly when they share an
endpoint. A graph $G$ is a line graph if it is isomorphic
to $L(H)$ for some finite simple graph $H$.

A graph is $4$-connected if it has at least five vertices
and remains connected after deletion of any set of at most
three vertices. A Hamilton cycle is a connected spanning
subgraph of degree two at every vertex, hence a simple
cycle visiting the entire graph.

The conjecture states that every $4$-connected line graph
has a Hamilton cycle. In terms of a root graph $H$, whenever
$L(H)$ is $4$-connected there should be a cyclic ordering
\[
 e_1,e_2,\ldots,e_m\qquad(m=|E(H)|)
\]
of all edges of $H$ such that $e_i$ and $e_{i+1}$ share
an endpoint for every index modulo $m$.

This ordering is the Hamilton cycle in $L(H)$. It need
not be an Euler tour in $H$: consecutively shared endpoints
need not give a consistent traversal of each original edge.
The hypothesis concerns vertex deletion in the line graph,
not vertex connectivity of $H$. In particular, it should
not be replaced by the different assumption that the root
graph itself is $4$-edge-connected. The root graph is useful
for expressing the target, but the conjecture is a property
of $G=L(H)$.
\subsection{Short English statement}
Does every line graph that remains connected after deletion of any three vertices contain a Hamilton cycle?
\subsection{Research attempt: an exact translation to a dominating Eulerian subgraph}
\paragraph{Scope, literature, and plan.}
GPT-6 Astra Ultra, 1 October 2026. Four-connectivity is imposed on
$L(H)$, not on the root $H$. Harary and Nash-Williams [S3] supply
the classical dominating Eulerian characterization; the argument is
reproduced below to expose the actual missing step. Kaiser and Vrána
[S4] prove Hamiltonicity for five-connected line graphs of minimum
degree at least six. That stronger hypothesis cannot simply replace
the present one, and no claim that it is the latest optimal partial
result is needed here.

\paragraph{1. A sufficient certificate in the root graph.}
Call a connected even-degree subgraph $D$ of $H$ dominating when every
edge of $H$ has at least one endpoint in $V(D)$. Allow $D$ to be
a single vertex without edges. Assume $|E(H)|\ge3$.
If $D$ has edges, follow an Euler tour of $D$, listing its edges
cyclically. For each edge outside $D$, choose one endpoint in $V(D)$
and assign the edge to it. At one visit to each such vertex $v$
in the tour, insert all edges assigned to $v$ between the incoming
and outgoing tour edges, in any order.

This cyclic list contains every edge of $H$ exactly once. Consecutive
edges share an endpoint: inside each inserted block they all meet
$v$, and the two adjacent tour edges also meet $v$. Thus it is
a Hamilton cycle in $L(H)$. If $D$ is a single vertex, domination
means every edge of $H$ meets that vertex, and any cyclic ordering
of those edges works. The certificate requires no spanning Euler tour
of the whole root graph.

\paragraph{2. Recovering that certificate from a line-graph cycle.}
Conversely, let $e_1,\ldots,e_m$ be a Hamilton cyclic ordering in
$L(H)$. For each $i$, let $v_i$ be the common endpoint of $e_i$
and $e_{i+1}$, with indices taken modulo $m$. The endpoint is unique
because $H$ is simple and the two edges are distinct. Both $v_{i-1}$
and $v_i$ are endpoints of $e_i$.

When $v_{i-1}\ne v_i$, traverse $e_i$ from the former to the latter.
When they agree, omit $e_i$ from this traversal and stay at the
shared vertex. Omitting these stationary steps leaves a closed trail
whose used edges are all distinct. If any edge is used, its edge
set $D$ is connected and has even degrees, because arrivals and
departures pair at every vertex. Every omitted edge meets a vertex
of the trail: along each stationary block the common vertex equals
the endpoint of its neighbouring nonstationary block. Thus $D$
is dominating. If no edge is used, all $v_i$ agree, and their one
common vertex is the degenerate dominating certificate. This proves
the equivalence in both directions.

\paragraph{3. Translating connectivity without strengthening it.}
Deleting at most three vertices of $L(H)$ corresponds to deleting
at most three edges of $H$. The remaining line graph is disconnected
exactly when the remaining root graph has at least two components
each containing an edge. Isolated root vertices produce no line-graph
vertices and must be ignored. Hence the hypothesis forbids small edge
separations into two edge-containing parts, rather than all small
edge cuts in the usual sense.

For example $H=K_{1,5}$ has bridges and vertices of odd degree, yet
$L(H)=K_5$ is four-connected and Hamiltonian. Its central vertex is
the degenerate dominating certificate. The script [S5] explicitly checks
the line-graph construction and its connectivity. This example rules
out replacing the desired certificate by a spanning even-degree
subgraph of the root: no such nontrivial spanning subgraph exists
in a star, though the original target holds.

\paragraph{Verification and first remaining gap.}
In the sufficiency proof every unused edge is assigned once, even
if both endpoints lie in $D$; no edge is duplicated by insertion.
In the converse, stationary steps are necessary: consecutive
line-graph adjacencies do not in general orient all root edges into
an Euler tour. The unresolved assertion is that the stated prohibition
on small edge separations always yields a dominating connected even
subgraph, possibly a single vertex. The equivalence identifies the
correct target but does not construct it from four-connectivity.

\subsection{Final status}
UNRESOLVED. The dominating Eulerian characterization is proved and the
connectivity translation is checked. No general construction at the
four-connected threshold is established. Completion estimate: no
defensible global percentage.

\subsection{Sources}
\begin{itemize}
\item[S1] ULAM.AI and the UnsolvedMath Contributors, supplied problems.json, record 3168 (OPG-485), OpenGarden collection. Catalogue identity and source transcription; CC BY 4.0.
\url{https://huggingface.co/datasets/ulamai/UnsolvedMath}.
\item[S2] Open Problem Garden, Hamiltonian cycles in line graphs. Original problem and discussion; the mathematical formulation below is expanded from the supplied source record.
\url{http://www.openproblemgarden.org/op/hamiltonian_cycles_in_line_graphs}.
\item[S3] F. Harary and C. St. J. A. Nash-Williams, \emph{On Eulerian and Hamiltonian Graphs and Line Graphs}, Canadian Mathematical Bulletin 8 (1965), 701--709. Checked 1 October 2026.
\url{https://doi.org/10.4153/CMB-1965-051-3}.
\item[S4] T. Kaiser and P. Vrána, \emph{Hamilton cycles in 5-connected line graphs}, arXiv:1009.3754, revised 2011. Checked 1 October 2026.
\url{https://arxiv.org/abs/1009.3754}.
\item[S5] GPT-6 Astra Ultra, exact line-graph example check, 1 October 2026: \texttt{scripts/check\_attempt\_batch02\_connectivity\_2026\_10\_01.py} in this repository; run with Python 3.
\end{itemize}
\end{document}
